As large language models (LLMs) ~\citep{openai2024gpt4technicalreport,dubey2024llama3herdmodels} demonstrate increasingly strong reasoning and planning capabilities, LLM-driven agentic systems have achieved remarkable performance in a wide range of tasks~\citep{Wang_2024}. 
Substantial effort has been invested in manually designing sophisticated agentic systems using human priors in different application areas. Recently, there has been a significant interest in creating self-evolving agents, that not only greatly reduce human labor but also produce better solutions. Given that human effort can only cover a small search space of agent design, it is reasonable to expect that a self-evolving agent with the freedom to explore the full design space has the potential to produce a more optimal solution.


\begin{figure}[t]
\setlength{\abovecaptionskip}{-0.1cm}
\setlength{\belowcaptionskip}{-0.4cm}
\begin{center}
%\framebox[4.0in]{$\;$}
% \fbox{\rule[-.5cm]{0cm}{4cm} \rule[-.5cm]{4cm}{0cm}}
\includegraphics[width=0.4\textwidth]{figures/intro.pdf}
\end{center}
\caption{Modular demonstration of \framework{}. Compared with traditional agents, its sensor and executor can read and write all of its own code. 
}
\label{fig:module}
\end{figure}

There is a large body of work proposing agents capable of self-refinement. Some agents are designed to iterate over a fixed routine consisting of a list of fixed modules, while some of the modules are capable of taking self- or environment feedback to refine their actions ~\citep{chen2023teachinglargelanguagemodels, qu2024tool, tang2025chemagentselfupdatinglibrarylarge}. This type of agent, referred to as \textbf{Hand-Designed Agent}, is depicted as having the lowest degree of freedom in \cref{fig:compare}. More automated agents have been designed to be able to update their routines or modules in some pre-defined meta-learning routine, for example, natural language gradients~\citep{zhou2024symbolic}, meta agent~\citep{hu2024automated}, or creating and collecting demonstrations ~\citep{khattab2023dspy}. This type of agent, known as \textbf{Meta-Learning Optimized Agents}, is depicted as having the middle degree of freedom in \cref{fig:compare}.
However, there are inevitably some human priors involved in these agent designs that cannot be improved during the inference time.




% such as reflection~\citep{shinn2024reflexion}, self-debugging~\citep{chen2023teachinglargelanguagemodels}, tool learning~\citep{qu2024tool}, and search algorithms~\citep{yao2023treethoughtsdeliberateproblem}. 
% \plm{[expand this part] Once these agents are designed, they typically operate within fixed frameworks, which constrains their ability to adapt and improve.} \yxj{mention figure 1}

% In response, some works have proposed using learning algorithms to optimize agents, such as accumulating experiences in external memory~\citep{zhang2024survey}, optimizing prompts based on environmental feedback~\citep{khattab2023dspy}, refining agent logic through natural language gradients~\citep{zhou2024symbolic}, or using LLMs to design new agents automatically~\citep{hu2024automated}. 
% While these learning algorithms can optimize agents, they are also based on human priors and are fixed, which can be a bottleneck in the improvement of agents. \plm{what are the bottlenecks? give more details / examples.} \yxj{mention figure 1}



% It is evident that both types of agents above are inherently constrained by human priors and one intuitional method to further increase the freedom of self-improvement is to design a meta-meta-learning algorithm, to learn the meta-learning algorithm. However, there is always a higher-level meta-learning algorithm that can be manually designed to learn the current-level meta-learning method, creating a never-ending hierarchy of meta-learning. 


\begin{figure*}[t]
\setlength{\abovecaptionskip}{-0.1cm}
\setlength{\belowcaptionskip}{-0.4cm}
\begin{center}
%\framebox[4.0in]{$\;$}
% \fbox{\rule[-.5cm]{0cm}{4cm} \rule[-.5cm]{4cm}{0cm}}
\includegraphics[width=0.9\textwidth]{figures/compare.pdf}
\end{center}
\caption{Comparison of three agent paradigms. Hand-designed agents rely on human expertise which are limited in scope and labor-intensive. Meta-learning optimized agents are constrained by a fixed meta-learning algorithm, restricting their search space and optimization potential. In contrast, self-referential agent (\framework) can \textbf{recursively} improve itself without any limitation. 
Its optimization capabilities are constantly being enhanced by itself. Consequently, in return, it can continue to optimize itself better.
% Note that the input to \framework is itself, allowing it to modify itself and output a new version of itself.
}
%\wxy{This is too hard to understand as a first figure. Please see comments in the figure compare\_xinyi.pdf. I agree with Liangming that a more concrete example is needed. Perhaps you can use a simple task as a running example, and solve it with 3 different level agents?}}
\label{fig:compare}
\end{figure*}
% \plm{Figure 1: a more concrete example is needed.} 


In this paper, we propose \textbf{\framework} to eliminate the human design prior, which is an automated LLM agent that can freely decide its own routine, modules, and even the way to update them. 
It is inspired by the self-referential \godel machine \citep{schmidhuber2003godel}, which was proven to be able to find the global optimal solutions. 
\textit{Self-reference} means the property of a system that can analyze and modify its own code, including the parts responsible for the analysis and modification processes~\citep{astrachan1994self}. Therefore, it can achieve what's known as "\textit{recursive self-improvement}", where it iteratively updates itself to become more efficient and effective at achieving its predefined goals. In this case, as shown in Figure \ref{fig:module}, \framework can analyze and modify its own code, including the code for analyzing and modifying itself, and thus can search the full agent design space, which is depicted as having the highest degree of freedom in \cref{fig:compare}. \framework can theoretically make increasingly better modifications over time through recursively self-update \citep{wang2018formulationrecursiveselfimprovementpossible}. 

% It benefits from the \textit{self-referential} nature of \framework. Self-reference means the property of a system that can analyze and modify its own code, including the parts responsible for the analysis and modification processes. Therefore, it can achieve what's known as "recursive self-improvement," where it iteratively updates itself to become more efficient and effective at achieving its predefined goals. This recursive nature allows it to theoretically make increasingly better modifications over time.
% Fixed meta-learning algorithms might not always be optimal due to their inability to adapt after initial deployment. This limitation necessitates the design of higher-order learning systems, such as meta-meta-learning, to continually refine and optimize the existing meta-learning frameworks. Theoretically, there's always a higher level of algorithm that can be designed to enhance the current meta-learning methods, introducing a potentially infinite regression of needs further optimization. In contrast, self-referential systems do not face such limitations; they inherently possess the ability to modify and improve themselves autonomously, bypassing the need for creating an endless hierarchy of learning algorithms to achieve optimal performance.



% Therefore, this raises the question: \plm{\textit{if the learning algorithms currently in use can be improved, do we need to design higher-level learning algorithms for them again? Then at this rate, this hierarchical structure will never converge}. ??? }
% Fortunately, such a hierarchical structure is not necessary. 
%To solve this dilemma, \citet{schmidhuber2003godel} proposed an optimal and self-referential \yxj{explain more about self-referential} problem solver named \godel machine that can optimize its own learning process without the need for higher-level learning algorithms. 
%The \godel machine can modify its initial learning algorithm and use the new method to continue optimizing all aspects in the manner of recursive self-improvement, unbounded by the limitations of fixed hierarchical designs~\citep{steunebrink2012towards}. 

%\plm{sounds incremental} Inspired by the \godel machine and leveraging the increasingly powerful LLMs, we propose a self-referential framework called \textbf{\framework}. This framework enables the agent to be dynamically aware of and modify its own logic, and at the next time step, the boosted one will have more power and recursively improve itself towards a given goal, leading to self-improvement. 
% Specifically, since most existing agents consist of code-based algorithms combined with decision-making capabilities derived from LLMs, if an agent can be aware of its own code\footnote{The code that \framework can read and modify includes the part related to its optimization logic.} and dynamically modify it (self-awareness and self-modification) based on its initial optimization principles, it can then continue this recursive process according to the newly updated optimization logic. The entire process is goal-driven and guided by the LLM's planning, allowing the agent to interact with the environment. Therefore, human intervention is no longer necessary, avoiding any design bottlenecks.
%\plm{These three different paradigms of agents are shown in Figure \ref{fig:compare}. [should talk about this part earlier.]}


In this paper, we choose to implement it by letting it manipulate its own runtime memory, i.e., the agent is able to retrieve its current code in the runtime memory and modify it by \textit{monkey patching} \citep{archiveMonkeyPatching}, which dynamically modifies classes or modules during execution.
% In our implementation, we adhere to a minimalist design to minimize the influence of human priors. 
% We implement the optimization module using a recursive function. 
To allow it to update the logic of the running main function, \textbf{unlike the loop-iterative approach of traditional agents, we implement the main function as a recursive function.}
In this function, LLM analyzes and makes a series of decisions, including reading and modifying its own code from runtime memory (\textit{self-awareness\footnote{In this paper, self-awareness means that the agent can introspect and read its own code and files, not to imply any philosophical sense of consciousness or awareness.}} and \textit{self-modification}), and interacting with the environment to gather feedback. The agent then proceeds to the subsequent recursive depth and continues to optimize itself. 
% It is worth noting that the optimization module may have already been modified by the time the recursion occurs, potentially enhancing its optimization capabilities.
%... \wxy{Need a more concrete description of your implementation here. Only mentioning the name of the modules is not enough.}
%\footnote{It is important to distinguish between program runtime memory and the external memory of the agent.}. 
% we refer to our implementation of the \framework as \textbf{\method}\footnote{Roz is a self-programming and self-evolving robot in the movie \textit{The Wild Robot}.}
% As shown in Figure \ref{fig:method}, \method dynamically retrieves its current code in the runtime memory and modifies it at training time. This approach is akin to a programming technique called “\textit{monkey patching}”, which allows for dynamic modification of classes or modules during execution. 
% In our initial implementation, we adhere to a minimalist design, minimizing the influence of human priors. 
% The key components of our implementation include thinking before acting, error handling, code runner and file processor.
% All of these initial implementations can be freely modified by \method during subsequent optimization processes as needed. %\yxj{should be more brief}

% \wxy{texts below are not fixed}  
% \wxy{How about let's just call both the framework and the implementation \framework? I don't see much benefit in calling them two separate names.}

% \plm{the advantage of your method is still unclear.}
To validate the effectiveness of \framework, we conduct experiments on multiple domains including coding, science, math, and reasoning. Our results demonstrate that \framework achieves significant performance gain across various tasks, surpassing various widely-used agents that require human design. The same implementation of \framework can easily adapt to different tasks by only specifying the environment description and feedback mechanism. Additionally, the case study of the optimization progress reveals that \framework can provide novel insights into agent design. Our codes are released to facilitate future research\footnote{\href{https://github.com/Arvid-pku/Godel\_Agent}{https://github.com/Arvid-pku/Godel\_Agent}}. 
% We also investigate the impact of the initial policy for improvement on subsequent outcomes, finding that a good start can significantly accelerate convergence during optimization.
% \wxy{define initial logic here}
% However, due to the current limitations of LLMs, \method occasionally falls into the trap of error accumulation, as noted by \citet{Yampolskiylimitation}, resulting in not being very robust. But we firmly believe that as LLM capabilities increase, recursive self-improving agents will completely outperform those elaborately hand-designed by humans.

% \plm{the concept "self-referential" is not well-explained.}

In summary, our contributions are as follows:
\vspace{-7pt}
\begin{itemize}[leftmargin=*,itemsep=1pt]
    \item We propose the first fully self-referential agent framework, \framework, and implement it using monkey patching. It autonomously engages in self-awareness, self-modification, and recursive self-improvement.
    % across any task, reducing the need for manual agent design and offering higher flexibility and freedom.  % \wxy{really?}. %This framework opens up an exciting new direction that warrants further exploration.
    \item Experiments shows that \framework{} is superior to the previous agent frameworks in terms of performance, flexibility, cost, and potential.
    % \item We implement \framework framework using the monkey patching. Experiments show that \framework outperforms manually designed agents, demonstrating effective self-improvement.
    \item We analyze \framework’s optimization process, including its self-referential abilities and the optimized agentic systems, aiming to deepen our understanding of both LLMs and agents.
    \item Our framework offers a promising direction for developing flexible and capable agents through recursive self-improvement.
\end{itemize}
